% R. Brown's translations of two letters to L. Breen, with his note and comments (cote 134-2): archive cote 134-2, site page 129 (PDF page 130).
% translation of the letter of 17-19 July 1975. Draft compared: ps1.tex, lines 1385-1394.
% AI-transcribed by Claude Opus 5.5 (claude-opus-5-5, adaptive thinking, effort high), Message Batches API, batch brown_full (msgbatch_01Md3i5yK1PxkpJn5GHfEExE), 2026-11-13; not yet checked by eye.
\scanpage{134-2/129}
\mspage{- 2 -}

(n+1)-equivalent to that of complexes of chains of length $n$ .
Agreed?  I thank you for having rectified in my mind a big blunder,
because of my great ignorance of algebraic topology and homotopy -
I was in fact of the impression that $H$-spaces satisfying conditions
of associativity and commutativity strict enough (say equivalent
to an $\Omega^{i}X$ with $i$ arbitrarily large) correspond to commutative
topological groups (and so inspiring in me several analogies ...).
Thus I have intuitive agreement with your observations on p.5.

On the other hand, I am always intrigued by the following
question: is there an analogue of the ``tapis'' of Dold-Puppe for
semi-simplicial groups (not necessarily commutative) and what form
should it take?  To say it truly I am considering the yoga
\[ \text{simplicial sets } \hand{$\longleftrightarrow$} \ \infty\text{-groupoids} \]
as being essentially the ultimate ``set theoretic'' version of
Dold-Puppe, which I would deduce by making explicit solely the
fact that the abelian groups in  ($\infty$-cat) are ``nothing else'' than
the complexes of chains in  (Ab) .  It is necessary then to
determine what should be the groups in ($\infty$-cat). I can tell you what
these are in (1-cat) - this will be discussed at length in the book
of Mme. Sinh, I think, under the title of the chapter ``\uline{strict}
Gr-categories'' (i.e. the isomorphisms of associativity, unity and
inverse  $X.X^{-1} \simeq 1$  are \uline{identities}).  One can make explicit how
(\uline{via} the fact that a Gr-category is Gr-equivalent to a strict
Gr-category)      the calculation with the Gr-categories reduces
to a very algebraic calculation with the \uline{strict} Gr-categories,
by a kind of ``calculus of fractions'' (by choice, left or right)
of the type which ends in giving the construction of derived

\mspage{\hand{130}}
